% =========================================================
% Proof: Inverse of a Bijective Isometry is an Isometry
% Source: notes/isometries/notes-isometries.tex
% Difficulty: Easy
% =========================================================

\subsection*{Inverse of a Bijective Isometry}
\label{prf:inverse-isometry}

\begin{remark*}[Return]
$\leftarrow$ \hyperref[prop:inverse-isometry]{Back to Proposition in Notes}
\end{remark*}

\begin{proof}
\Claim If $f : X \to Y$ is a bijective isometry then $f^{-1} : Y \to X$
is a bijective isometry.

\Given A bijective isometry $f : X \to Y$.

\Goal To show $d_X(f^{-1}(u), f^{-1}(v)) = d_Y(u,v)$ for all $u,v \in Y$.

\Strategy Write $u = f(x)$ and $v = f(y)$ (possible since $f$ is bijective).
Then $f^{-1}(u) = x$, $f^{-1}(v) = y$, and $d_X(x,y) = d_Y(f(x),f(y)) = d_Y(u,v)$.

% -------------------------------------------------------
% YOUR PROOF HERE
% -------------------------------------------------------

\end{proof}

\begin{remark*}[Proof shape]
Change variables using surjectivity, then apply the original isometry condition.
\end{remark*}

\begin{remark*}[Dependencies]
\begin{itemize}
  \item Definition~\ref{def:isometry}: $d_Y(f(x),f(y)) = d_X(x,y)$.
  \item Bijectivity: $f$ is both injective and surjective.
\end{itemize}
\end{remark*}
